%% ============================================================
%% 分片 34：原文第 11 节「$A$-Branes and ${\cal D}$-Modules」（原文 glshort.tex 起于第 9938 行）
%% 由 tools/harvmac2tex.py 从 arXiv 源码（hep-th/0604151v3）自动生成。
%% 只译英文散文；公式、编号、\label、\cite 一律不改。
%% ============================================================

\subsection{Generalizations of the c.c. Brane and Twisted Differential
Operators}\label{sec:genotics}

We have seen that geometric Langlands duality can be understood as
a result of the $S$-transformation exchanging $\Psi=0$ and
$\Psi=\infty$. More generally, the transformation $S$ maps a gauge
theory with gauge group $G$ and canonical parameter $\Psi$ to the
gauge theory with gauge group $\LG$ and the canonical parameter
$-1/n_{\frak g}\Psi$ (where, as we explained in section 第~\ref{sec:reviewsd}~节,
$n_{\frak g}=1$ if $G$ is simply-laced, and otherwise $n_{\frak
g}=2$ or 3). Our goal here is to describe what sort of
generalization of the geometric Langlands program arises when we
depart from $\Psi=0,\infty$. (See \cite{ref9}, sections 6.3 and 8.6,
for a discussion of this generalization via conformal field
theory, and \cite{ref146} for the abelian case.)

The generalization of the geometric Langlands program that we will
explore is perhaps most interesting for the case of rational
$\Psi$, since, as we saw in section 第~\ref{sec:tophooft}~节, this is the case in
which there are interesting line operators. Moreover, the full
$S$-duality group, which as we showed in section 第~\ref{sec:canonpar}~节, is
generated by $T:\Psi\to\Psi+1$ along with $S:\Psi\to -1/n_{\frak
g}\Psi$, maps rational values of $\Psi$ to rational values. In
fact, all rational values of $\Psi$ are related by this group, a
fact which also ensures that at a rational value of $\Psi$ there
are plenty of branes, just as there are at $\Psi=0$ and
$\Psi=\infty$.

The duality nonetheless works for all complex-valued $\Psi$. For
simplicity, however, here we will take $\Psi$ to be real. One can
obtain an arbitrary real value of $\Psi$ by letting $t=1$ and
keeping the $\theta$-angle arbitrary: % 式 (11.26)
\setcounter{equation}{25}
\begin{equation}\label{eq:melo}
\Psi={\theta\over
{2\pi}}=\Re\tau.
\end{equation}
 $S$-duality in general transforms ${\rm
Re}\,\tau=\theta$, ${\rm Im}\,\tau=4\pi/e^2$, and $t$. But $\Psi$
is the only important parameter, so as long as $\Psi$ is real, we
can think of $S$-duality as acting  on $\theta$ only.  In this
set-up, the ``electric'' and ``magnetic'' theories look much more
symmetric: they both have $t=1$ but their $\theta$-angles are
inversely related. Upon compactification on the Riemann surface
$C$, one obtains an $A$-model in complex structure $K$ with target
$\MH(G,C)$ and $B$-field % 式 (11.27)
\setcounter{equation}{26}
\begin{equation}\label{eq:pelbo}
B=-\omega_I \Re\tau,
\end{equation}
 as
explained in eqn. \eqref{eq:palooka}. The $S$-transformation is implemented
by replacing $\Re\tau$ by $-1/n_{\frak g}\Re\tau$ and $G$ by
$^L\negthinspace G$.

The $B$-field given by \eqref{eq:pelbo}\ is not flat along the fibers of the
Hitchin fibration; indeed, it is a multiple of the Kahler form in
complex structure $I$ and hence is nondegenerate when restricted
to a fiber. An important consequence is that the fibers of the
Hitchin fibration are no longer valid $A$-branes for general
$\Psi$. To explain this, recall that the sigma-model with target
$X$ in the presence of a $D$-brane wrapped on a submanifold $Y$ is
invariant under a gauge transformation % 式 (11.28)
\setcounter{equation}{27}
\begin{equation}\label{eq:zebol}
B\to
B-d\lambda,\quad A\to A+\lambda\vert_Y,\quad \lambda\in
\Omega^1(X).
\end{equation}
 The gauge-invariant combination is $F+B\vert_Y$. Now
suppose $X$ is symplectic and $Y$ is Lagrangian. If $B=0$, then a
single $D$-brane wrapped on $Y$ can be an $A$-brane only if $F=0$.
The  gauge-invariant generalization of this condition, which is
also valid for $B\neq 0$, is % 式 (11.29)
\setcounter{equation}{28}
\begin{equation}\label{eq:subol}
F+B\vert_Y=0.
\end{equation}
 If the
cohomology class of $B|_Y/2\pi$ is not integral, then this
equation cannot be satisfied, for any $U(1)$ gauge field, and
therefore a single $A$-brane cannot be wrapped on $Y$. In our
case, this means that for $\theta\neq 2\pi n$, the fibers of the
Hitchin fibration are not valid $A$-branes, regardless of the
choice of Chan-Paton bundle.

For non-integer $\Psi$, the objects of study will therefore not be
branes of rank one supported on a fiber of the Hitchin fibration.
 As a consistency check, note that their naive $T$-duals, i.e.
points on $\MH(\LG,C)$, are also not valid topological branes at
$t=1$ and $\theta\neq 0$.  In fact, a point is never an $A$-brane,
since the support of an $A$-brane is always at least
middle-dimensional.  For rational $\Psi$, there are $A$-branes of
finite rank supported on a fiber of the Hitchin fibration; we can
find what they are by following the duality starting at $\Psi=0$.
In any case, $S$-duality still maps  $A$-branes for $\MH(G,C)$ to
$A$-branes for $\MH(\LG,C)$, even though examples of $A$-branes
may be scarce.

But what can we say about $A$-branes when $\Psi\not=0$ that
generalizes what we have found at $\Psi=0$? In the remainder of
this subsection we will answer this question.
 We will perform all computations near the classical limit $e^2\to
0,\ \Im\,\tau\to \infty$, but since the theory depends only on
$\Psi=\Re\tau$, the results will be valid for all values of $e^2$.

We will see that an analogue of the c.c. brane exists for all
$\Psi$.  We will call this brane ${\cal B}_{c.c.}^\Psi$. Its
algebra of open strings will turn out to be ${\cal
D}_{K_{\EUM}^{1/2}\otimes {{{\frak L}}}^\Psi}$, that is, the sheaf
of holomorphic differential operators on $K_\M^{1/2}\otimes{\frak
L}^\Psi$, where ${\frak L}$ is the determinant line bundle on
$\cM(G,C)$ (see the discussion of eqn. \eqref{eq:frox}). Given any other
$A$-brane ${\cal B}'$ on $\MH(G,C)$, the sheaf of $({\cal
B}_{c.c.}^\Psi,{\cal B}')$ strings is a module for ${\cal D}_{K_{
\EUM}^{1/2}\otimes {{{\frak L}}}^\Psi}$.  So an $A$-brane at
general $\Psi$ gives a sheaf of modules for this algebra. Just as
at $\Psi=0$, the construction really comes from four-dimensional
gauge theory, which in algebraic geometry means that these twisted
${\cal D}$-modules can be extended to the stack of all
$G$-bundles. $S$-duality then maps  twisted ${\cal D}$-modules for
$G$ to those of  $^L\negthinspace G$, with the exchange $\Psi\to
-1/n_{\frak g}\Psi$.

The gauge field on a space-filling  $A$-brane must solve the
equation % 式 (11.30)
\setcounter{equation}{29}
\begin{equation}\label{eq:juko}
(\omega^{-1}(F+B))^2=-1,
\end{equation}
 where
$\omega=(\Im\tau)\, \omega_K$. An obvious solution is % 式 (11.31)
\setcounter{equation}{30}
\begin{equation}\label{eq:uko}
F=
\Im\tau\, (\cos{\eurm q})\, \omega_J,\quad \sin{\eurm q}=-{{\Re\,
\tau}\over {\Im\, \tau}}.
\end{equation}
  We take this solution to define the
brane ${\cal B}_{c.c.}^\Psi$.  Note that, as ${\Im\,\tau}$ can be
arbitrarily large, ${\eurm q}$ is extremely small, and ultimately
the interesting effects will be linear in $\Psi=\Re\,\tau$, which
is of order ${\eurm q}$.

If we ignore the algebra structure, then the sheaf of open strings
is locally isomorphic to the sheaf of holomorphic functions on
$\MH(G,C)$ in complex structure\footnote{Comparing to eqn. \eqref{eq:wyre}, we
see that $I(\Psi)$ is the same as $I_w$ with $w=i\tan({\eurm
q}/2)$.} % 式 (11.32)
\setcounter{equation}{31}
\begin{equation}\label{eq:muko}
I(\Psi)=\omega^{-1}(F+B)=I \cos{\eurm q} -
J\sin{\eurm q}.
\end{equation}
 The corresponding holomorphic coordinates on
$\MH(G,C)$ are % 式 (11.33)
\setcounter{equation}{32}
\begin{equation}\label{eq:suko}
A'_\bz=A_\bz - i \tan{{\eurm q}\over
2}\, \phi_\bz,\quad \phi'_z=\phi_z+i\tan{{\eurm q}\over 2}\, A_z.
\end{equation}

As usual, $A$ is a connection on a $G$-bundle $E\to C$, and $\phi$
is a one-form valued in ${\rm ad}(E)$. What we have in section
第~\ref{sec:hitchfib}~节\ called Hitchin's second fibration is the forgetful map
% 式 (11.34)
\setcounter{equation}{33}
\begin{equation}\label{eq:xnonko}
\pi_0: (A_\bz,\phi_z)\to A_\bz
\end{equation}
 from $\MH$ to $\M$. It
is plain that $\pi_0$ is not holomorphic with respect to $I(\Psi)$
except for $\Psi=0$. To rectify the situation, let us define a new
$\bpartial$-connection on $E$ by % 式 (11.35)
\setcounter{equation}{34}
\begin{equation}\label{eq:hoxo}
D'_\bz=D_\bz-i
\tan{{\eurm q}\over 2}\,\phi_\bz=\partial_\bz+A_\bz - i
\tan{{\eurm q}\over 2}\,\phi_\bz,
\end{equation}
 and map the pair $(A,\phi)$
into the holomorphic bundle on $C$ defined by $D'_\bz$. Let us
denote this map $\pi_\Psi$; it is obviously holomorphic with
respect to $I(\Psi)$ and for $\Psi=0$ reduces to $\pi_0$. As
usual, we will restrict this fibration to the subset of Higgs
bundles for which $E$, endowed with the modified $\bar\partial$
operator, is stable.

What does the fiber of the map $\pi_\Psi$ look like? To understand
this, we note that the obvious candidate for the fiber coordinate,
namely $\phi_z+i\tan{{\eurm q}\over 2}\, A_z$, transforms
inhomogeneously under gauge transformations. The object which
transforms homogeneously is % 式 (11.36)
\setcounter{equation}{35}
\begin{equation}\label{eq:gyko}
D'_z=\phi_z+i\tan{{\eurm
q}\over 2}\, D_z,
\end{equation}
 where $D_z$ is the holomorphic covariant
derivative on $E$. Up to a factor $i\tan{{\eurm q}\over 2}$,
$D'_z$ is a $\partial$-connection over the bundle $E$. We choose
not to divide by $i\tan{{\eurm q}\over 2}$, because we would like
the formulas to make sense even for $\Psi=0$. Then $D'_z$ is what
is called a holomorphic $\lambda$-connection on $E$, where
$\lambda=i\tan{{\eurm q}\over 2}$. For any complex manifold $X$, a
holomorphic $\lambda$-connection on a holomorphic bundle $E\to X$
is a linear map $\nabla: \Gamma(E)\ra \Gamma(E\ot T^*X)$ that
commutes with the $\bar\partial$ operator of $E$ and such that
% 式 (11.37)
\setcounter{equation}{36}
\begin{equation}\label{eq:hyfo}
\nabla (f\cdot s)=f\nabla s + \lambda
\partial f\wedge s,
\end{equation}
 for any function $f$ and any  section $s$ of
$E$. For $\lambda=1$, a $\lambda$-connection is an ordinary
$\partial$ operator on $E$, while for $\lambda=0$ it is a section
of $\End(E)\ot T^*X$.

By virtue of Hitchin's equations, $D'_z$ commutes with $D'_\bz$
(this is closely related to the discussion of complex structure
$I_w$ at the end of section 第~\ref{sec:complexstr}~节) and therefore defines a
holomorphic $\lambda$-connection on the holomorphic bundle
$E'=\pi_\Psi(E,\phi)$. By the usual logic, this implies that the
fiber over $E'$ is the moduli space of holomorphic
$\lambda$-connections on $E'$.

The space of holomorphic $\lambda$-connections on a fixed vector
bundle $E'\to C$ is naturally an affine space modeled on the space
of holomorphic sections of $\End(E')\ot K_C$, where $K_C$ is the
canonical line bundle of $C$. (This statement means that to a
holomorphic $\lambda$-connection, we can add a holomorphic section
of $\End(E')\otimes K_C$, and any two holomorphic
$\lambda$-connections are related in this way.)  The space
$H^0(C,\End(E')\ot K_C)$ can be identified with the cotangent
space to $\cM(G,C)$ at the point $E'$. Thus in complex structure
$I(\Psi)$, the moduli space $\MH(G,C)$ looks like an affine
deformation of $T^*\cM(G,C)$; in other words, it is a bundle of
affine spaces, with the underlying bundle of vector spaces being
$T^*\CM(G,C)$. To summarize, if we ignore the algebra structure,
the sheaf of $({\cal B}_{c.c.}^\Psi,{\cal B}_{c.c.}^\Psi)$ strings
is the sheaf of holomorphic functions on a ``twisted'' cotangent
bundle of $\cM(G,C)$.

It remains to determine the algebra structure of the sheaf of
$({\cal B}_{c.c}^\Psi,{\cal B}_{c.c.}^\Psi)$ strings. We restrict
ourselves to functions of polynomial growth along the fibers of
the affine bundle.\footnote{This restriction amounts to picking a
particular algebraic structure on  $\MH$ endowed with the complex
structure $I(\Psi)$. In this algebraic structure, it is called the
moduli space of $G$-bundles with $\lambda$-connection. There is a
second algebraic structure that is compatible with the same
complex structure, namely one in which $\MH$ is equivalent to the
moduli space of stable homomorphisms $\pi_1(C)\to G_{\Bbb{C}}$. In
the second algebraic structure, the traces of holonomies are
regarded as the algebraic functions. What we have just encountered
is the one point in the present paper in which we have to choose
between the different algebraic structures compatible with a fixed
complex structure on $\MH$.  This choice plays a more prominent
role in other approaches to the geometric Langlands program.} We
can follow the same logic as at $\Psi=0$.
 Locally, the fact that we are quantizing a twisted cotangent
bundle rather than an ordinary one is irrelevant. Globally, it
might (and as we will see, it does) modify the transition
functions in \eqref{eq:hecox}. So the sheaf of c.c. strings at any $\Psi$ is
the sheaf of holomorphic differential operators acting on a power
of some holomorphic line bundle over $\cM(G,C)$.

We can be more precise if $\Psi$ is an integer. Shifting $\Psi$ to
$\Psi+n$ shifts the $B$-field, or more precisely $B/2\pi$, by the
first Chern class of the line bundle ${{{\frak L}}}^{-n}$. (We can
see this from \eqref{eq:palooka}; the shift $\Psi\to \Psi+n$, which is
$\theta\to\theta+2\pi n$, shifts $B$ by $-n\omega_I$, where
$\omega_I/2\pi$ is the first Chern class of ${\frak L}$.) This
observation enables us to define a convenient $A$-brane   in
complex structure $K$. At $\Psi=0$, one of the important branes
was the brane ${\cal B}_0$ supported at $\phi=0$ with trivial
Chan-Paton bundle. At $\Psi=n$, we likewise can consider a brane
${\cal B}_n$ supported at $\phi=0$, but now with Chan-Paton bundle
${{{\frak L}}}^n$, endowed with its natural connection.  This
obeys \eqref{eq:subol}, and so is an $A$-brane in complex structure $K$.

Since we have set $\Psi=n$, we write ${\cal B}_{c.c.}^n$ instead
of ${\cal B}_{c.c.}^\Psi$. Now we quantize the $({\cal
B}_{c.c.}^n,{\cal B}_n)$ strings. Everything is as before, except
that we must include the Chan-Paton bundle of the brane ${\cal
B}_n$.  So the physical states are now sections of
$K_{\EUM}^{1/2}\otimes {\frak
L}^n\otimes\left(\oplus_{j=0}^n\Omega^{0,j}(\M)\right)$.  This is
a module for ${\cal D}_{K_{\EUM}^{1/2}\otimes {{{\frak L}}}^n}$,
and not for ${\cal D}_{K_{\EUM}^{1/2}\otimes {{{\frak L}}}^m}$ for
any other $m$. So the sheaf of algebras of $({\cal
B}_{c.c.}^n,{\cal B}_{c.c.}^n)$ strings is precisely ${\cal
D}_{K_{\EUM}^{1/2}\otimes {{{\frak L}}}^n}$.

What happens if $\Psi$ is not an integer?  We cannot hope to find
for general $\Psi$ a brane that leads to a module consisting of
sections of $K_{\EUM}^{1/2}\otimes {{{\frak L}}}^\Psi$, since
there is no such module unless $\Psi$ is an integer. However, we
know that the algebra of c.c. strings at any $\Psi$ is the sheaf
of differential operators acting on ${K_{\EUM}^{1/2}}\otimes
{\frak L}^{f(\Psi)}$ for some function $f(\Psi)$.  So if we can
prove that $f(\Psi)$ is a linear function, this, together with the
fact that $f(\Psi)=\Psi$ for integer $\Psi$, implies that
$f(\Psi)=\Psi$ for all $\Psi$.

To show that $f(\Psi)$ is linear, consider the Euclidean action of
the sigma-model: % 式 (11.38)
\setcounter{equation}{37}
\begin{equation}\label{eq:cumber}
S=\int_\Sigma\Phi^*
(\omega-iF-iB)=-i\,\Im\tau \int_\Sigma \Phi^*\left(\cos{\eurm q}\,
\omega_J+\sin{\eurm q}\, \omega_I + i\omega_K\right).
\end{equation}
 (We have
omitted terms of the form $\{Q,V\}$.)
 For $\Psi=0$, we have ${\eurm q}=0$, and the two-form in
the integrand becomes an exact two-form on $T^*\cM(G,C)$. The
corresponding one-form potential $\varpi$ is, up to a numerical
factor, the canonical one-form $p_\alpha\, dq^\alpha$.  For
$\Psi\neq 0$, the best we can do is to write this two-form as an
exact form on $T^*\cM(G,C)$ plus the pull-back of a closed
two-form from $\cM(G,C)$: % 式 (11.39)
\setcounter{equation}{38}
\begin{equation}\label{eq:hexo}
\Omega'={{\Im\tau}\over
\pi}\delta\int_C \,\Tr\left(\phi_z \delta A'_\bz\right)-
i\,\Re\tau \int_C\, \Tr\left(\delta A'_z\delta A'_\bz\right),
\end{equation}

where we defined $A'_z$ as minus the Hermitian conjugate of
$A'_\bz$. The second term in this formula is $\Re\tau$ times a
$(1,1)$-form  on $\M(G,C)$ which is closed but not exact.  It is a
multiple of the curvature of the natural line bundle ${\frak
L}\to\M$. To compute the transition functions \eqref{eq:hecox}, we do
perturbation theory using a propagator that is the inverse of the
form $\Omega'$. \eqref{eq:hecox}\ was deduced from the commutator $[\hat
p_\alpha,\hat p_\beta]$, which is proportional to $\hbar^2$ or
${1/ ({\rm Im}\,\tau)^2}$. So we need to do perturbation theory up
to order $1/({\rm Im}\,\tau)^2$. In that order, perturbation
theory is linear in $\Psi={\rm Re}\,\tau$, since the inverse of
$\Omega'$ can be expanded as a series in ${\rm Re}\,\tau$ and
${\rm Im}\,\tau$ that schematically is of the form $\sum_{k\geq
0}a_k({\rm Re}\,\tau)^k/({\rm Im}\,\tau)^{k+1}$. This completes
the argument.

The canonical line bundle over $\cM(G,C)$ is isomorphic to ${\frak
L}^{-2h}$, where $h$ is the dual Coxeter number of $G$
\cite{ref141} . Thus we can also say that the algebra of c.c.
strings at any given value of $\Psi$ is the sheaf of holomorphic
differential operators on the line bundle ${\frak L}^{\Psi-h}$. In
particular, if we let $\Psi =h$, the algebra of c.c. strings is
the sheaf of ordinary differential operators, that is,
differential operators acting on functions. It follows that by
combining $S$-duality with a transformation by $T^h$, which shifts
$\Psi$ by $h$, we can map a zerobrane on $\MH(\LG,C)$ to an
ordinary (untwisted) ${\cal D}$-module on $\cM(G,C)$. This is a
more standard formulation of the geometric Langlands duality. For
certain $G$, it is also a more precise statement, as one sees if
one considers carefully \cite{ref5,ref26} the dependence on the
spin structure of $C$.
